\chapter{无约束最优化方法}\label{chap:unconstrained}
% 第1节：引言
\section{引言}\label{sec:ch5-intro}

\subsection{为什么研究无约束最优化}

实际优化问题通常含有大量约束，但研究无约束最优化方法具有重要意义：

\begin{enumerate}
    \item \textbf{转化工具}：许多算法可通过Lagrangian乘子法将有约束问题转化为一系列无约束问题，如Lagrangian对偶、鞍点最优性条件、惩罚函数法和障碍函数法。
    
    \item \textbf{基础组件}：大多数方法通过寻找方向并沿方向最小化来推进，这种线搜索本质上是无约束或简单约束的最小化问题。
    
    \item \textbf{自然推广}：几种无约束最优化方法可自然推广到有约束问题。
\end{enumerate}

\subsection{问题描述}

无约束最优化问题的标准形式：
\[
\min_{x \in \mathbb{R}^n} f(x)
\]

\begin{conceptbox}[关键假设]
\begin{itemize}
    \item \textbf{零阶方法}：仅需函数值 $f(x)$
    \item \textbf{一阶方法}：需要梯度 $\nabla f(x)$
    \item \textbf{二阶方法}：需要Hessian矩阵 $\nabla^2 f(x)$
\end{itemize}
\end{conceptbox}

\subsection{方法分类}

本笔记涵盖的方法体系如下：

\begin{notebox}[关于方法分类]
本章聚焦\textbf{多维无约束最优化}的核心算法。一维线搜索方法（如二分搜索、黄金分割法、Fibonacci法）和非精确线搜索准则（Wolfe条件、Armijo规则）已在第~\ref{chap:linesearch}~章系统介绍，作为多维优化算法中的线搜索子程序。本章涵盖所有多维无约束优化方法，从仅需函数值的零阶方法到需要导数和Hessian的高阶方法。
\end{notebox}


\begin{table}[H]
\centering
\caption{无约束最优化方法分类}
\begin{tabular}{@{}lll@{}}
\toprule
\textbf{类别} & \textbf{方法} & \textbf{所需信息} \\
\midrule
\multirow{3}{*}{多维无导数}
 & 坐标轮换法 & 函数值 \\
 & Hooke-Jeeves方法 & 函数值 \\
 & Rosenbrock方法 & 函数值 \\
\midrule
\multirow{2}{*}{一阶方法} 
 & 最速下降法（梯度法） & 梯度 \\
 & 共轭梯度法 & 梯度 \\
\midrule
\multirow{3}{*}{二阶方法} 
 & 牛顿法 & Hessian \\
 & Levenberg-Marquardt & Hessian修正 \\
 & 信任域方法 & Hessian近似 \\
\midrule
拟牛顿法 & DFP、BFGS & 梯度 \\
\midrule
非光滑优化 & 次梯度方法 & 次梯度 \\
\bottomrule
\end{tabular}
\end{table}

\subsection{收敛速度记号}

\begin{definition}[收敛速度]
设序列 $\{x_k\}$ 收敛到 $x^*$：
\begin{itemize}
    \item \textbf{线性收敛}：$\|x_{k+1} - x^*\| \leq q \|x_k - x^*\|$，$q \in (0,1)$
    \item \textbf{超线性收敛}：$\|x_{k+1} - x^*\| / \|x_k - x^*\| \to 0$
    \item \textbf{二次收敛}：$\|x_{k+1} - x^*\| \leq C \|x_k - x^*\|^2$
\end{itemize}
\end{definition}

% 第2节：最速下降法

% ==================== 从第四章移入 ====================
% 多维无导数方法本质上属于n维无约束优化，
% 从第四章移至第五章，使第五章覆盖全部多维无约束方法。
% ====================

\section{多维无导数搜索方法}\label{sec:derivative-free}

对于问题 $\min_{x \in \mathbb{R}^n} f(x)$，当导数信息不可得或计算代价高昂时，使用仅依赖函数值的方法。

\subsection{坐标轮换法（Cyclic Coordinate Method）}\label{sec:cyclic}

\begin{conceptbox}[基本思想]
沿坐标轴方向依次搜索，每次只改变一个变量。搜索方向为标准基向量 $e_1, e_2, \ldots, e_n$。
\end{conceptbox}

\textbf{算法步骤：}
\begin{algorithm}[H]
\caption{坐标轮换法}
\begin{algorithmic}[1]
\Require 初始点 $x^{(0)}$，精度 $\varepsilon > 0$
\State $k \gets 0$
\While{未收敛}
    \State $y^{(0)} \gets x^{(k)}$
    \For{$i = 1$ \textbf{to} $n$}
        \State 求 $\lambda_i = \arg\min f(y^{(i-1)} + \lambda e_i)$
        \State $y^{(i)} = y^{(i-1)} + \lambda_i e_i$
    \EndFor
    \State $x^{(k+1)} \gets y^{(n)}$
    \If{$\|x^{(k+1)} - x^{(k)}\| < \varepsilon$}
        \State \textbf{break}
    \EndIf
    \State $k \gets k + 1$
\EndWhile
\State \Return $x^{(k+1)}$
\end{algorithmic}
\end{algorithm}

\begin{example}[坐标轮换法]
求解 $\min f(x) = x_1^2 - x_1 x_2 + x_2^2 - 2x_1$，初始点 $(0, 3)$，最优解 $(2, 1)$。

\textbf{解：}

\textbf{第1轮迭代}：

沿 $e_1 = (1, 0)$ 搜索：$\min f(\lambda, 3) = \lambda^2 - 3\lambda + 9 - 2\lambda$
\[
\frac{df}{d\lambda} = 2\lambda - 5 = 0 \Rightarrow \lambda = 2.5
\]
得 $y^{(1)} = (2.5, 3)$

沿 $e_2 = (0, 1)$ 搜索：$\min f(2.5, 3 + \lambda) = 6.25 - 2.5(3+\lambda) + (3+\lambda)^2 - 5$
\[
\frac{df}{d\lambda} = -2.5 + 6 + 2\lambda = 0 \Rightarrow \lambda = -1.75
\]
得 $x^{(1)} = (2.5, 1.25)$

继续迭代直至收敛。
\end{example}

\begin{notebox}[坐标轮换法的特点]
\begin{itemize}
    \item 每轮迭代进行 $n$ 次一维搜索
    \item 函数可微时收敛到梯度为零的点
    \item 不可微时可能在非最优点停止
    \item 收敛速度与最速下降法相当
    \item 可使用Gauss-Southwell变体：每次选偏导数绝对值最大的方向
\end{itemize}
\end{notebox}

\subsection{Hooke-Jeeves方法}\label{sec:hooke-jeeves}

\begin{conceptbox}[基本思想]
结合两种搜索策略：
\begin{itemize}
    \item \textbf{试探搜索（Exploratory Search）}：沿各坐标方向试探，寻找改进方向
    \item \textbf{模式搜索（Pattern Search）}：沿成功方向加速前进
\end{itemize}
\end{conceptbox}

\textbf{算法步骤：}
\begin{enumerate}
    \item \textbf{初始化}：$x^{(0)}$（基点），步长 $\delta > 0$，缩减因子 $\rho \in (0, 1)$
    
    \item \textbf{Step 1（试探搜索）}：
    从 $y = x^{(k)}$ 出发，沿各坐标方向 $e_i$ 试探：
    \begin{itemize}
        \item 若 $f(y + \delta e_i) < f(y)$，则 $y \leftarrow y + \delta e_i$（成功）
        \item 否则若 $f(y - \delta e_i) < f(y)$，则 $y \leftarrow y - \delta e_i$
        \item 否则保持 $y$ 不变（失败）
    \end{itemize}
    设结果为 $x^{(k+1)}$
    
    \item \textbf{Step 2（模式搜索判断）}：
    \begin{itemize}
        \item 若 $f(x^{(k+1)}) < f(x^{(k)})$（试探成功）：
        \begin{itemize}
            \item 沿模式方向 $d = x^{(k+1)} - x^{(k)}$ 前进：$y = x^{(k+1)} + d = 2x^{(k+1)} - x^{(k)}$
            \item $x^{(k)} \leftarrow x^{(k+1)}$，从 $y$ 转Step 1
        \end{itemize}
        \item 否则（试探失败）：
        \begin{itemize}
            \item 若 $\delta < \varepsilon$，停止
            \item 否则 $\delta \leftarrow \rho \delta$，从 $x^{(k)}$ 转Step 1
        \end{itemize}
    \end{itemize}
\end{enumerate}

\begin{example}[Hooke-Jeeves方法]
求解 $\min f(x) = x_1^2 - x_1 x_2 + x_2^2 - 2x_1$，初始点 $(0, 0)$，$\delta = 0.5$。

\textbf{解：}

\textbf{第1轮}：

从 $(0,0)$ 试探：
- $e_1$ 方向：$f(0.5, 0) = -0.75 < f(0,0) = 0$ (成功)，$y = (0.5, 0)$
- $e_2$ 方向：$f(0.5, 0.5) = -0.5 < f(0.5, 0) = -0.75$ (失败)，$f(0.5, -0.5) = 0.25 > -0.75$ (失败)

得 $x^{(1)} = (0.5, 0)$，试探成功。

模式方向 $d = (0.5, 0) - (0, 0) = (0.5, 0)$，新模式基点：
$y = (0.5, 0) + (0.5, 0) = (1, 0)$

继续迭代直至满足精度要求。
\end{example}

\subsection{Rosenbrock方法}\label{sec:rosenbrock}

\begin{conceptbox}[基本思想]
每次迭代沿 $n$ 个线性无关的正交方向进行一维搜索，搜索完成后根据搜索轨迹构造一组新的正交方向（Gram-Schmidt正交化）。
\end{conceptbox}

\textbf{算法步骤：}

\textbf{Step 1}：给定初始点 $x^{(0)}$，初始正交方向 $d_1, \ldots, d_n$（通常取坐标轴），步长 $\delta_1, \ldots, \delta_n$

\textbf{Step 2}：沿各方向依次搜索：
\begin{itemize}
    \item 对 $i = 1, \ldots, n$：
    \begin{itemize}
        \item 若 $f(x + \delta_i d_i) < f(x)$，成功：$x \leftarrow x + \delta_i d_i$，$\delta_i \leftarrow \alpha \delta_i$（$\alpha > 1$）
        \item 否则失败：$\delta_i \leftarrow -\beta \delta_i$（$\beta \in (0, 1)$）
    \end{itemize}
\end{itemize}

\textbf{Step 3}：一轮完成后若所有方向都有成功搜索，转Step 2

\textbf{Step 4}：若存在未成功的方向：
\begin{itemize}
    \item 若 $|\delta_i| < \varepsilon$（所有 $i$），停止
    \item 否则，利用Gram-Schmidt正交化构造新的正交方向，转Step 2
\end{itemize}

\begin{notebox}[Gram-Schmidt正交化]
设第 $k$ 轮搜索的位移为 $\Delta x_1, \ldots, \Delta x_n$，新方向构造为：
\begin{align*}
s_1 &= \Delta x_1, \quad d_1^{new} = \frac{s_1}{\|s_1\|} \\
s_i &= \Delta x_i - \sum_{j=1}^{i-1} \frac{\Delta x_i^T s_j}{s_j^T s_j} s_j, \quad d_i^{new} = \frac{s_i}{\|s_i\|}
\end{align*}
\end{notebox}

\subsection{Python代码实现}\label{sec:df-code}

\begin{codebox}[坐标轮换法]
\begin{lstlisting}
import numpy as np
from scipy.optimize import minimize_scalar

def cyclic_coordinate(f, x0, eps=1e-6, max_iter=1000):
    """
    坐标轮换法
    
    Parameters:
        f: 目标函数
        x0: 初始点
        eps: 精度
        max_iter: 最大迭代次数
    
    Returns:
        x_opt: 最优解
        f_opt: 最优值
        history: 迭代历史
    """
    n = len(x0)
    x = x0.copy()
    history = [x.copy()]
    
    for _ in range(max_iter):
        x_old = x.copy()
        
        for i in range(n):
            # 沿第i个坐标方向搜索
            def phi(lam):
                x_new = x.copy()
                x_new[i] += lam
                return f(x_new)
            
            res = minimize_scalar(phi)
            x[i] += res.x
        
        history.append(x.copy())
        
        if np.linalg.norm(x - x_old) < eps:
            break
    
    return x, f(x), history

# 示例
def f(x):
    return x[0]**2 - x[0]*x[1] + x[1]**2 - 2*x[0]

x0 = np.array([0.0, 3.0])
x_opt, f_opt, hist = cyclic_coordinate(f, x0)
print(f"最优解: {x_opt}")
print(f"最优值: {f_opt:.6f}")
\end{lstlisting}
\end{codebox}

\begin{codebox}[Hooke-Jeeves方法]
\begin{lstlisting}
def hooke_jeeves(f, x0, delta=0.5, rho=0.5, eps=1e-6, max_iter=1000):
    """
    Hooke-Jeeves模式搜索法
    
    Parameters:
        f: 目标函数
        x0: 初始点
        delta: 初始步长
        rho: 步长缩减因子
        eps: 精度
        max_iter: 最大迭代次数
    
    Returns:
        x_opt: 最优解
        history: 迭代历史
    """
    n = len(x0)
    x_base = x0.copy()
    x = x0.copy()
    history = [x.copy()]
    
    for _ in range(max_iter):
        # 试探搜索
        y = x_base.copy()
        success = False
        
        for i in range(n):
            e_i = np.zeros(n)
            e_i[i] = 1
            
            # 正向试探
            if f(y + delta * e_i) < f(y):
                y = y + delta * e_i
                success = True
            # 负向试探
            elif f(y - delta * e_i) < f(y):
                y = y - delta * e_i
                success = True
        
        if success:
            # 模式搜索
            x = y + (y - x_base)
            x_base = y.copy()
            history.append(y.copy())
        else:
            if delta < eps:
                break
            delta *= rho
            x = x_base.copy()
    
    return x_base, history

# 示例
x0 = np.array([0.0, 0.0])
x_opt, hist = hooke_jeeves(f, x0, delta=0.5)
print(f"最优解: {x_opt}")
\end{lstlisting}
\end{codebox}

\subsection{方法比较}

\begin{table}[H]
\centering
\caption{多维无导数方法比较}
\begin{tabular}{@{}llll@{}}
\toprule
\textbf{方法} & \textbf{搜索方向} & \textbf{步长策略} & \textbf{收敛性} \\
\midrule
坐标轮换 & 固定坐标轴 & 精确线搜索 & 线性收敛 \\
Hooke-Jeeves & 坐标轴+模式方向 & 离散步长 & 局部收敛 \\
Rosenbrock & 正交方向（自适应） & 离散步长 & 局部收敛 \\
\bottomrule
\end{tabular}
\end{table}

\section{最速下降法（梯度法）}\label{sec:gradient}

\subsection{基本思想}

\begin{conceptbox}[最速下降方向]
1847年Cauchy提出。在点 $x$ 处，归一化的最速下降方向为：
\[
d = -\frac{\nabla f(x)}{\|\nabla f(x)\|}
\]
即负梯度方向，是在单位球面上使函数局部下降最快的方向。
\end{conceptbox}

\begin{lemma}[最速下降方向的验证]
设 $f: \mathbb{R}^n \to \mathbb{R}$ 在 $x$ 处可微，$\nabla f(x) \neq 0$。求解：
\[
\min_{\|d\| \leq 1} \nabla f(x)^T d
\]
其解为 $d = -\frac{\nabla f(x)}{\|\nabla f(x)\|}$。
\end{lemma}

\textbf{证明}：由Cauchy-Schwarz不等式：
\[
\nabla f(x)^T d \geq -\|\nabla f(x)\| \cdot \|d\| \geq -\|\nabla f(x)\|
\]
等号成立当且仅当 $d = -\frac{\nabla f(x)}{\|\nabla f(x)\|}$。

\subsection{最速下降法算法}

\begin{algorithm}[H]
\caption{最速下降法}
\begin{algorithmic}[1]
\Require 初始点 $x^{(0)}$，精度 $\varepsilon > 0$
\State $k \gets 0$
\While{$\|\nabla f(x^{(k)})\| \geq \varepsilon$}
    \State $d_k = -\nabla f(x^{(k)})$
    \State 求 $\lambda_k = \arg\min_{\lambda \geq 0} f(x^{(k)} + \lambda d_k)$
    \State $x^{(k+1)} = x^{(k)} + \lambda_k d_k$
    \State $k \gets k + 1$
\EndWhile
\State \Return $x^{(k)}$
\end{algorithmic}
\end{algorithm}

\begin{example}[最速下降法]
求解 $\min f(x) = x_1^2 + 4x_2^2$，初始点 $\begin{bmatrix} 4 \\ 1 \end{bmatrix}$。

\textbf{解：}

$\nabla f(x) = (2x_1, 8x_2)$，$x^{(0)} = (4, 1)$

\textbf{第1次迭代}：$d_0 = (-8, -8)$

$\phi(\lambda) = f(4-8\lambda, 1-8\lambda) = (4-8\lambda)^2 + 4(1-8\lambda)^2$

$\phi'(\lambda) = -16(4-8\lambda) - 64(1-8\lambda) = -128 + 640\lambda = 0$

解得 $\lambda_0 = 0.2$，$x^{(1)} = (2.4, -0.6)$

\textbf{第2次迭代}：$d_1 = (-4.8, 4.8)$

$\phi(\lambda) = f(2.4-4.8\lambda, -0.6+4.8\lambda) = (2.4-4.8\lambda)^2 + 4(-0.6+4.8\lambda)^2$

$\phi'(\lambda) = -46.08 + 230.4\lambda = 0$

解得 $\lambda_1 = 0.2$，$x^{(2)} = (1.44, 0.36)$，继续迭代。
\end{example}

\subsection{步长选择策略}\label{sec:step-size}

\subsubsection{策略1：固定步长}

$\lambda_k = \lambda$（常数）。

\begin{theorem}[固定步长收敛条件]
设 $f \in C^{1,1}(\mathbb{R}^n)$（Lipschitz连续梯度，常数为 $L$），若 $0 < \lambda < \frac{2}{L}$，则：
\[
f(x^{(k+1)}) - f(x^*) \leq f(x^{(k)}) - f(x^*) - \lambda\left(1 - \frac{L\lambda}{2}\right)\|\nabla f(x^{(k)})\|^2
\]
\end{theorem}

最优固定步长：$\lambda = \frac{1}{L}$。

\subsubsection{策略2：精确线搜索（全松弛）}

$\lambda_k = \arg\min_{\lambda \geq 0} f(x^{(k)} + \lambda d_k)$。

\begin{theorem}[精确线搜索的收敛率]
设 $f$ 为强凸函数（参数 $\mu$）且 $L$-光滑，则：
\[
f(x^{(k+1)}) - f^* \leq \left(\frac{L - \mu}{L + \mu}\right)^2 (f(x^{(k)}) - f^*)
\]
即线性收敛，收敛因子为 $\frac{L-\mu}{L+\mu} = 1 - \frac{2}{L/\mu + 1}$。
\end{theorem}

\begin{notebox}[条件数的影响]
定义条件数 $\kappa = \frac{L}{\mu}$（Hessian最大与最小特征值之比）。

收敛因子 $\approx 1 - \frac{2}{\kappa}$，当 $\kappa$ 很大时（病态问题），收敛极慢。
\end{notebox}

\subsubsection{策略3：Armijo规则}

实际中最常用的策略（详见第\ref{sec:armijo}节）。

\subsection{梯度法的收敛性分析}\label{sec:gradient-convergence}

\begin{theorem}[梯度下降下降量]
设 $f \in C^{1,1}$ 且 $L$-光滑，采用步长 $\lambda_k = \frac{1}{L}$，则：
\[
f(x^{(k+1)}) \leq f(x^{(k)}) - \frac{1}{2L}\|\nabla f(x^{(k)})\|^2
\]
\end{theorem}

\textbf{证明}：由Lipschitz梯度的二次上界：
\[
f(y) \leq f(x) + \nabla f(x)^T(y-x) + \frac{L}{2}\|y-x\|^2
\]
代入 $y = x - \frac{1}{L}\nabla f(x)$：
\[
f(y) \leq f(x) - \frac{1}{L}\|\nabla f(x)\|^2 + \frac{1}{2L}\|\nabla f(x)\|^2 = f(x) - \frac{1}{2L}\|\nabla f(x)\|^2
\]

\subsubsection{凸函数情形}

\begin{theorem}[凸函数收敛率]
设 $f$ 为凸函数且 $L$-光滑，$\|x^{(0)} - x^*\| \leq R$，则梯度法（步长 $\frac{1}{L}$）满足：
\[
f(x^{(k)}) - f^* \leq \frac{2LR^2}{k}
\]
即收敛速率为 $O(1/k)$。
\end{theorem}

\subsubsection{强凸函数情形}

\begin{theorem}[强凸函数线性收敛]
设 $f$ 为 $\mu$-强凸且 $L$-光滑，则：
\[
\|x^{(k)} - x^*\|^2 \leq \left(1 - \frac{\mu}{L}\right)^k \|x^{(0)} - x^*\|^2
\]
即线性收敛。
\end{theorem}

\subsection{梯度法的局限性}\label{sec:gradient-limitation}

\begin{enumerate}
    \item \textbf{锯齿现象}：相邻搜索方向正交，在最优解附近路径呈锯齿状
    \item \textbf{条件数敏感}：对病态问题（$\kappa \gg 1$）收敛极慢
    \item \textbf{只能保证收敛到稳定点}：不一定收敛到局部极小点
    \item \textbf{鞍点困难}：可能停滞在鞍点
\end{enumerate}

\begin{figure}[H]
\centering
\begin{tikzpicture}[scale=0.72]
    % 坐标轴
    \draw[->,thick] (-0.3,0) -- (7.5,0) node[right] {$x_1$};
    \draw[->,thick] (0,-0.3) -- (0,5.5) node[above] {$x_2$};
    % 等值线（椭圆，条件数大——病态）
    \foreach \a/\b in {0.5/1.5, 1/3, 1.5/4.5, 2/6} {
        \draw[blue!30] (3,2.5) ellipse ({\a} and {\b*0.5});
    }
    % 最优点
    \fill[red] (3,2.5) circle (3pt) node[below right,scale=0.8] {$x^*$};
    % 锯齿路径
    \draw[->,very thick,red!70!black] (0.5,0.5) -- (2.5,0.5);
    \draw[->,very thick,red!70!black] (2.5,0.5) -- (2.5,4);
    \draw[->,very thick,red!70!black] (2.5,4) -- (3.8,4);
    \draw[->,very thick,red!70!black] (3.8,4) -- (3.8,1.5);
    \draw[->,very thick,red!70!black] (3.8,1.5) -- (3.2,1.5);
    \draw[->,very thick,red!70!black] (3.2,1.5) -- (3.2,3);
    \draw[->,very thick,red!70!black] (3.2,3) -- (3.5,3);
    % 起点
    \fill (0.5,0.5) circle (2.5pt) node[below left,scale=0.75] {$x^{(0)}$};
    % 标注
    \node[scale=0.8,red!70!black] at (1,4.8) {锯齿路径};
    % 条件数标注
    \node[scale=0.7,blue!60] at (6,4.5) {\small 椭圆越扁};
    \node[scale=0.7,blue!60] at (6,4) {\small 条件数越大};
    \node[scale=0.7,blue!60] at (6,3.5) {\small 收敛越慢};
\end{tikzpicture}
\caption{最速下降法的锯齿现象：对病态问题（大条件数）收敛缓慢}
\label{fig:gradient-zigzag}
\end{figure}

\begin{example}[非凸函数的梯度法]
$f(x, y) = x^3 - 3x + y^3 - 3y$

\textbf{解：}

稳定点：$(1,1)$（极小），$(-1,-1)$（极大），$(1,-1)$ 和 $(-1,1)$（鞍点）

以 $(-1, 0.5)$ 为初始点，梯度法可能收敛到鞍点 $(-1, 1)$。
\end{example}

\subsection{梯度法的Python实现}\label{sec:gradient-code}

\begin{codebox}[最速下降法完整实现]
\begin{lstlisting}
import numpy as np
from scipy.optimize import minimize_scalar

def gradient_descent(f, grad, x0, method='exact', alpha=1e-4,
                     rho=0.5, L=None, mu=None, eps=1e-6, 
                     max_iter=1000):
    """
    最速下降法
    
    Parameters:
        f: 目标函数
        grad: 梯度函数
        x0: 初始点
        method: 'exact'(精确线搜索), 'fixed'(固定步长), 
                'armijo'(Armijo回溯)
        alpha: Armijo参数
        rho: 步长缩减因子
        L: Lipschitz常数（固定步长时使用）
        eps: 梯度范数精度
        max_iter: 最大迭代次数
    
    Returns:
        x_opt: 最优解
        f_opt: 最优值
        history: 迭代历史
    """
    x = x0.copy()
    history = [x.copy()]
    
    for k in range(max_iter):
        g = grad(x)
        
        if np.linalg.norm(g) < eps:
            break
        
        d = -g  # 最速下降方向
        
        if method == 'exact':
            # 精确线搜索
            def phi(lam):
                return f(x + lam * d)
            res = minimize_scalar(phi, bounds=(0, 10))
            lam = res.x
        elif method == 'fixed':
            # 固定步长
            lam = 1.0 / L if L else 0.01
        elif method == 'armijo':
            # Armijo回溯
            lam = 1.0
            fx = f(x)
            grad_dot_d = np.dot(g, d)
            for _ in range(20):
                if f(x + lam * d) <= fx + alpha * lam * grad_dot_d:
                    break
                lam *= rho
        
        x = x + lam * d
        history.append(x.copy())
    
    return x, f(x), history

# 示例：二次函数
def f(x):
    return x[0]**2 + 4*x[1]**2

def grad(x):
    return np.array([2*x[0], 8*x[1]])

x0 = np.array([4.0, 1.0])
x_opt, f_opt, hist = gradient_descent(f, grad, x0, method='exact')
print(f"最优解: {x_opt}")
print(f"最优值: {f_opt:.8f}")
print(f"迭代次数: {len(hist)}")
\end{lstlisting}
\end{codebox}

\subsection{批量与小批量梯度下降}\label{sec:sgd}

\begin{conceptbox}[随机梯度下降（SGD）]
在机器学习中，目标函数为经验风险：
\[
f(x) = \frac{1}{N}\sum_{i=1}^{N} f_i(x)
\]

\textbf{批量梯度下降（Batch GD）}：每次用全部数据
\[
x_{k+1} = x_k - \eta \nabla f(x_k) = x_k - \frac{\eta}{N}\sum_{i=1}^{N} \nabla f_i(x_k)
\]

\textbf{随机梯度下降（SGD）}：每次用一个样本
\[
x_{k+1} = x_k - \eta \nabla f_{i_k}(x_k)
\]
其中 $i_k$ 随机选取。

\textbf{小批量梯度下降（Mini-Batch GD）}：每次用 $m$ 个样本
\[
x_{k+1} = x_k - \frac{\eta}{m}\sum_{i \in B_k} \nabla f_i(x_k)
\]
\end{conceptbox}

\begin{notebox}[SGD的特点]
\begin{itemize}
    \item 计算效率高，适合大规模数据
    \item 不保证单调下降
    \item 学习率选择困难
    \item 鞍点问题严重
    \item 需要学习率衰减策略
\end{itemize}
\end{notebox}

% 第6节：牛顿法
\section{牛顿法}\label{sec:newton}

\subsection{基本思想}

\begin{conceptbox}[牛顿法]
利用二阶Taylor展开在 $x_k$ 处近似：
\[
f(x) \approx f(x_k) + \nabla f(x_k)^T(x-x_k) + \frac{1}{2}(x-x_k)^T \nabla^2 f(x_k)(x-x_k)
\]

令梯度为0，得迭代公式：
\[
x_{k+1} = x_k - [\nabla^2 f(x_k)]^{-1} \nabla f(x_k)
\]
即搜索方向 $d_k = -H_k^{-1} g_k$（$H_k$ 为Hessian，$g_k$ 为梯度）。
\end{conceptbox}

\begin{figure}[H]
\centering
\begin{tikzpicture}[scale=0.8]
    % 坐标轴
    \draw[->,thick] (-0.3,0) -- (6.5,0) node[right] {$x$};
    \draw[->,thick] (0,-0.3) -- (0,4.5) node[above] {$f(x)$};
    % 原函数
    \draw[very thick,blue]
        (0.3,4) .. controls (1,2) and (2,0.5) .. (3,0.3) .. controls (4,0.5) and (5,2) .. (6,4);
    \node[blue,scale=0.8] at (6.2,4.2) {$f$};
    % 当前点 x_k
    \fill[red] (1,2) circle (2.5pt) node[below left,scale=0.75] {\small $x_k$};
    % 二次近似（抛物线）
    \draw[red,thick,dashed]
        (0.3,3.5) .. controls (1.5,1.2) and (2.5,0.8) .. (3,1.5) .. controls (3.5,2.2) .. (4.5,4);
    \node[red,scale=0.7] at (4.8,4.2) {\small 二次近似};
    % 牛顿步
    \draw[->,very thick,green!50!black] (1,2) -- (2.8,0.35);
    \node[green!50!black,scale=0.75,above] at (1.6,1.6) {\small 牛顿步};
    % 新点 x_{k+1}
    \fill[green!50!black] (2.8,0.35) circle (2.5pt) node[below,scale=0.75] {\small $x_{k+1}$};
    % 最优解
    \fill[black] (3,0.3) circle (2pt) node[below,scale=0.7] {\small $x^*$};
    % 切线（梯度方向）
    \draw[orange,dotted,thick] (0.3,3) -- (2.5,0);
    \node[orange,scale=0.65] at (0.3,3.2) {\small 切线};
    % 标注
    \node[scale=0.7,gray] at (5,1) {\small 二阶收敛};
    \node[scale=0.7,gray] at (5,0.6) {\small 步收敛};
\end{tikzpicture}
\caption{牛顿法：用二次Taylor展开的极小点作为下一次迭代点}
\label{fig:newton-method}
\end{figure}

\subsection{牛顿法算法}

\begin{algorithm}[H]
\caption{牛顿法}
\begin{algorithmic}[1]
\Require 初始点 $x^{(0)}$，精度 $\varepsilon > 0$
\State $k \gets 0$
\While{$\|\nabla f(x^{(k)})\| \geq \varepsilon$}
    \State 计算 $g_k = \nabla f(x^{(k)})$
    \State 计算 $H_k = \nabla^2 f(x^{(k)})$
    \State 解线性方程组 $H_k d_k = -g_k$
    \State $x^{(k+1)} = x^{(k)} + d_k$
    \State $k \gets k + 1$
\EndWhile
\State \Return $x^{(k)}$
\end{algorithmic}
\end{algorithm}

\begin{example}[牛顿法解二次函数]
$f(x) = x_1^2 + 4x_2^2$，初始点 $\begin{bmatrix} 4 \\ 1 \end{bmatrix}$。

\textbf{解：}

$g(x) = \begin{bmatrix} 2x_1 \\ 8x_2 \end{bmatrix}$，
$H(x) = \begin{bmatrix} 2 & 0 \\ 0 & 8 \end{bmatrix}$

$H^{-1} = \begin{bmatrix} 0.5 & 0 \\ 0 & 0.125 \end{bmatrix}$

$d_0 = -H^{-1} g(4,1) = -\begin{bmatrix} 0.5 & 0 \\ 0 & 0.125 \end{bmatrix}\begin{bmatrix} 8 \\ 8 \end{bmatrix} = \begin{bmatrix} -4 \\ -1 \end{bmatrix}$

$x^{(1)} = (4,1) + (-4,-1) = (0, 0) = x^*$！

对二次函数，牛顿法一步收敛。
\end{example}

\subsection{牛顿法的收敛性分析}\label{sec:newton-convergence}

\begin{theorem}[局部二次收敛]
设 $f \in C^2(\mathbb{R}^n)$，$\nabla^2 f(x^*)$ 正定，$\nabla^2 f$ 在 $x^*$ 附近Lipschitz连续。则存在 $\delta > 0$，使得当 $\|x^{(0)} - x^*\| < \delta$ 时：
\[
\|x^{(k+1)} - x^*\| \leq C \|x^{(k)} - x^*\|^2
\]
即二次收敛。
\end{theorem}

\textbf{证明概要}：
\begin{align*}
x_{k+1} - x^* &= x_k - x^* - H_k^{-1} g_k \\
&= H_k^{-1}(H_k(x_k - x^*) - g_k) \\
&= H_k^{-1}(\nabla^2 f(x_k)(x_k - x^*) - (\nabla f(x_k) - \nabla f(x^*)))
\end{align*}
由 $\nabla^2 f$ 的Lipschitz连续性：
\[
\|\nabla f(x_k) - \nabla f(x^*) - \nabla^2 f(x_k)(x_k - x^*)\| = O(\|x_k - x^*\|^2)
\]

\subsection{牛顿法的优缺点}

\begin{table}[H]
\centering
\caption{牛顿法的优缺点}
\begin{tabular}{@{}ll@{}}
\toprule
\textbf{优点} & \textbf{缺点} \\
\midrule
局部二次收敛 & 需要计算Hessian矩阵 \\
对二次函数一步收敛 & 需解线性方程组 $O(n^3)$ \\
尺度不变性 & Hessian可能不正定（非凸问题） \\
& 初始点需充分接近最优解 \\
\bottomrule
\end{tabular}
\end{table}

\subsection{牛顿法的改进策略}\label{sec:newton-improvements}

\subsubsection{阻尼牛顿法（Damped Newton）}

引入步长参数：
\[
x_{k+1} = x_k - \lambda_k H_k^{-1} g_k
\]
其中 $\lambda_k$ 由线搜索确定（如Armijo条件）。

\subsubsection{修正牛顿法（Modified Newton）}

当 $H_k$ 不正定时，修正为：
\[
\tilde{H}_k = H_k + \nu_k I
\]
其中 $\nu_k > 0$ 使得 $\tilde{H}_k$ 正定（如取 $\nu_k = \max(0, \delta - \lambda_{\min}(H_k))$）。

\subsubsection{Cholesky分解修正}

对 $H_k$ 进行Cholesky分解时，若遇到非正定增广对角线元素，强制其为正数。

\subsection{牛顿法的Python实现}\label{sec:newton-code}

\begin{codebox}[牛顿法完整实现]
\begin{lstlisting}
import numpy as np
from scipy.linalg import solve

def newton_method(f, grad, hess, x0, line_search=True,
                  alpha=1e-4, rho=0.5, eps=1e-6, max_iter=100):
    """
    牛顿法
    
    Parameters:
        f: 目标函数
        grad: 梯度函数
        hess: Hessian矩阵函数
        x0: 初始点
        line_search: 是否使用线搜索
        eps: 精度
        max_iter: 最大迭代次数
    
    Returns:
        x_opt: 最优解
        f_opt: 最优值
        history: 迭代历史
    """
    x = x0.copy()
    history = [x.copy()]
    
    for k in range(max_iter):
        g = grad(x)
        
        if np.linalg.norm(g) < eps:
            break
        
        H = hess(x)
        
        # 解线性方程组 H * d = -g
        try:
            d = solve(H, -g)
        except np.linalg.LinAlgError:
            # Hessian奇异，回退到最速下降
            d = -g
        
        if line_search:
            # Armijo回溯线搜索
            lam = 1.0
            fx = f(x)
            grad_dot_d = np.dot(g, d)
            if grad_dot_d >= 0:  # 非下降方向
                d = -g
                grad_dot_d = np.dot(g, d)
            for _ in range(20):
                if f(x + lam * d) <= fx + alpha * lam * grad_dot_d:
                    break
                lam *= rho
        else:
            lam = 1.0
        
        x = x + lam * d
        history.append(x.copy())
    
    return x, f(x), history

# 示例: Rosenbrock函数
def rosenbrock(x):
    return (1 - x[0])**2 + 100 * (x[1] - x[0]**2)**2

def rosen_grad(x):
    return np.array([
        -2*(1 - x[0]) - 400*x[0]*(x[1] - x[0]**2),
        200*(x[1] - x[0]**2)
    ])

def rosen_hess(x):
    return np.array([
        [2 + 1200*x[0]**2 - 400*x[1], -400*x[0]],
        [-400*x[0], 200]
    ])

x0 = np.array([-1.0, 2.0])
x_opt, f_opt, hist = newton_method(rosenbrock, rosen_grad, 
                                   rosen_hess, x0)
print(f"最优解: {x_opt}")
print(f"最优值: {f_opt:.10f}")
print(f"迭代次数: {len(hist)}")
\end{lstlisting}
\end{codebox}

\begin{codebox}[带Hessian修正的牛顿法]
\begin{lstlisting}
def modified_newton(f, grad, hess, x0, eps=1e-6, max_iter=100):
    """
    修正牛顿法（处理非正定Hessian）
    """
    x = x0.copy()
    history = [x.copy()]
    
    for k in range(max_iter):
        g = grad(x)
        
        if np.linalg.norm(g) < eps:
            break
        
        H = hess(x)
        
        # 特征值修正
        eigvals, eigvecs = np.linalg.eigh(H)
        eigvals = np.maximum(eigvals, 1e-4)  # 强制正
        
        # H^{-1} = Q * Lambda^{-1} * Q^T
        H_inv = eigvecs @ np.diag(1.0 / eigvals) @ eigvecs.T
        d = -H_inv @ g
        
        # Armijo线搜索
        lam = 1.0
        fx = f(x)
        for _ in range(20):
            if f(x + lam * d) <= fx + 1e-4 * lam * np.dot(g, d):
                break
            lam *= 0.5
        
        x = x + lam * d
        history.append(x.copy())
    
    return x, f(x), history
\end{lstlisting}
\end{codebox}

\begin{notebox}[最速下降法、牛顿法、拟牛顿法的统一形式]
上述方法都可统一写成
\[
x_{k+1} = x_k + \alpha_k d_k, \quad d_k = -B_k^{-1} g_k, \quad g_k = \nabla f(x_k)
\]
不同方法对应不同的 $B_k$：
\begin{center}
\begin{tabular}{@{}lll@{}}
\toprule
\textbf{方法} & $B_k$ & \textbf{搜索方向} \\
\midrule
最速下降法 & $I$ & $d_k = -g_k$ \\
牛顿法 & $\nabla^2 f(x_k)$ & $d_k = -[\nabla^2 f(x_k)]^{-1} g_k$ \\
修正牛顿法 & $\nabla^2 f(x_k) + E_k$（保证正定） & 保证下降方向 \\
拟牛顿法 & Hessian 近似 $B_k$ 或逆近似 $H_k$ & $d_k = -H_k g_k$ \\
\bottomrule
\end{tabular}
\end{center}
拟牛顿法的核心思想是：不直接计算 Hessian，而是通过梯度差 $y_k = g_{k+1} - g_k$ 和步长差 $s_k = x_{k+1} - x_k$ 构造满足割线方程 $B_{k+1} s_k = y_k$（或 $H_{k+1} y_k = s_k$）的矩阵更新。DFP 和 BFGS 是典型代表。
\end{notebox}

% 第7节：Levenberg-Marquardt方法与信任域方法
\section{牛顿方法的变种：LM与信任域方法}\label{sec:lm-trust}

\subsection{Levenberg-Marquardt（LM）方法}\label{sec:lm}

\begin{conceptbox}[LM方法]
LM方法是处理非线性最小二乘问题的牛顿型方法，通过自适应调节在梯度下降和牛顿法之间切换：

对于问题：$\min f(x) = \frac{1}{2}\sum_{i=1}^{m} r_i(x)^2 = \frac{1}{2}\|r(x)\|^2$

迭代公式：
\[
(J_k^T J_k + \lambda_k I) d_k = -J_k^T r(x_k)
\]
其中 $J(x)$ 是残差向量 $r(x)$ 的Jacobian矩阵。
\end{conceptbox}

\textbf{LM方法的解释}：
\begin{itemize}
    \item 当 $\lambda_k \to 0$：趋近于高斯-牛顿法（$J^T J \approx$ Hessian的近似）
    \item 当 $\lambda_k \to \infty$：趋近于梯度下降（$d_k \approx -\frac{1}{\lambda_k} J^T r$）
\end{itemize}

\subsubsection{LM算法}

\begin{algorithm}[H]
\caption{Levenberg-Marquardt方法}
\begin{algorithmic}[1]
\Require 初始点 $x_0$，初始参数 $\lambda_0 > 0$，$\nu > 1$
\State $k \gets 0$
\While{未收敛}
    \State 计算 $J_k = J(x_k)$，$r_k = r(x_k)$
    \State 解 $(J_k^T J_k + \lambda_k I) d_k = -J_k^T r_k$
    \State 计算实际下降与预测下降比：
    \State $\rho_k = \dfrac{f(x_k) - f(x_k + d_k)}{f(x_k) - q_k(d_k)}$
    \State 其中 $q_k(d) = f(x_k) + (J_k^T r_k)^T d + \frac{1}{2}d^T(J_k^T J_k + \lambda_k I)d$
    \If{$\rho_k > 0$}
        \State $x_{k+1} = x_k + d_k$
        \State 若 $\rho_k > 0.75$，$\lambda_{k+1} = \lambda_k / \nu$
        \State 否则若 $\rho_k < 0.25$，$\lambda_{k+1} = \nu \lambda_k$
    \Else
        \State $x_{k+1} = x_k$
        \State $\lambda_{k+1} = \nu \lambda_k$
    \EndIf
    \State $k \gets k + 1$
\EndWhile
\end{algorithmic}
\end{algorithm}

\begin{example}[曲线拟合]
拟合模型 $y = ae^{bx}$ 到数据点 $(t_i, y_i)$：

\textbf{解：}

$\min f(a,b) = \sum_{i=1}^{n} (ae^{bt_i} - y_i)^2$

残差：$r_i(a,b) = ae^{bt_i} - y_i$

Jacobian：$J_{i1} = e^{bt_i}$，$J_{i2} = ate^{bt_i}$

用LM方法迭代求解。
\end{example}

\subsection{信任域方法（Trust Region Methods）}\label{sec:trust-region}

\begin{conceptbox}[信任域方法]
与线搜索不同，信任域方法先确定搜索范围（信任域），再在此范围内优化模型：
\[
\min_d q_k(d) = f(x_k) + g_k^T d + \frac{1}{2}d^T H_k d, \quad \text{s.t. } \|d\| \leq \Delta_k
\]
其中 $\Delta_k > 0$ 为信任域半径。
\end{conceptbox}

\subsubsection{信任域方法的框架}

\begin{enumerate}
    \item \textbf{求解子问题}：在给定半径 $\Delta_k$ 内求 $d_k$
    \item \textbf{评估比率}：
    \[
    \rho_k = \frac{f(x_k) - f(x_k + d_k)}{q_k(0) - q_k(d_k)}
    \]
    \item \textbf{更新信任域半径}：
    \begin{itemize}
        \item 若 $\rho_k$ 接近1：增大 $\Delta_k$
        \item 若 $\rho_k$ 较小：减小 $\Delta_k$
        \item 若 $\rho_k > \eta_0$：接受 $x_{k+1} = x_k + d_k$
        \item 否则：拒绝步，$x_{k+1} = x_k$
    \end{itemize}
\end{enumerate}

\subsubsection{信任域子问题的求解}

\begin{theorem}[子问题最优性条件]
$d^*$ 是信任域子问题的全局最优解当且仅当存在 $\lambda \geq 0$ 使得：
\begin{align}
(H + \lambda I)d^* &= -g \\
\lambda(\Delta - \|d^*\|) &= 0 \\
H + \lambda I &\succeq 0
\end{align}
\end{theorem}

\textbf{Cauchy点}：
在信任域内沿最速下降方向的最优点：
\[
d_k^C = -\tau_k \frac{\Delta_k}{\|g_k\|} g_k
\]
其中：
\[
\tau_k = \begin{cases} 1 & \text{if } g_k^T H_k g_k \leq 0 \\ \min\left(\frac{\|g_k\|^3}{\Delta_k g_k^T H_k g_k}, 1\right) & \text{otherwise} \end{cases}
\]

\subsubsection{Dogleg方法}

\begin{conceptbox}[Dogleg方法]
结合最速下降方向和牛顿方向的折线路径：
\begin{enumerate}
    \item 若牛顿步在信任域内（$\|d^N\| \leq \Delta$），取 $d = d^N$
    \item 否则，沿最速下降到Cauchy点的路径
    \item 若Cauchy点在信任域外（$\|d^{SD}\| \geq \Delta$），取边界点
    \item 否则在折线路径上取信任域边界点
\end{enumerate}
\end{conceptbox}

\begin{example}[信任域方法——一次迭代]\label{ex:trust-region}
用信任域方法求解
\[
\min\ f(x_1, x_2) = x_1^2 + x_2^2 - 3x_1 x_2 + 4x_1 - 5x_2 + 2
\]
初始点 $x_0 = (0, 0)$，信任域半径 $\Delta_0 = 1$，执行一次迭代。

\textbf{解：}

\textbf{Step 1：计算梯度与Hessian}
\[
\nabla f(x) = \begin{bmatrix} 2x_1 - 3x_2 + 4 \\ 2x_2 - 3x_1 - 5 \end{bmatrix},\quad
H = \begin{bmatrix} 2 & -3 \\ -3 & 2 \end{bmatrix}
\]

在 $x_0 = (0, 0)$ 处：$g_0 = (4, -5)^T$，$f(x_0) = 2$。

\textbf{Step 2：求Cauchy点}
\[
\|g_0\|^2 = 41,\quad g_0^T H g_0 = 2(16) - 6(4)(-5) + 2(25) = 202 > 0
\]
\[
\tau = \frac{\|g_0\|^3}{\Delta_0 \cdot g_0^T H g_0} = \frac{41\sqrt{41}}{202} = \frac{262.5}{202} \approx 1.30 > 1
\]
取 $\tau = 1$，Cauchy步沿最速下降方向走到信任域边界：
\[
d^C = -\frac{\Delta_0}{\|g_0\|} g_0 = -\frac{1}{\sqrt{41}} \begin{bmatrix} 4 \\ -5 \end{bmatrix} \approx \begin{bmatrix} -0.625 \\ 0.781 \end{bmatrix}
\]

\textbf{Step 3：计算模型下降量}
\[
g_0^T d^C = -\frac{41}{\sqrt{41}} = -\sqrt{41} \approx -6.403
\]
\[
(d^C)^T H d^C = \frac{202}{41} \approx 4.927
\]
\[
q_0(d^C) = g_0^T d^C + \frac{1}{2}(d^C)^T H d^C = -6.403 + 2.463 = -3.940
\]

\textbf{Step 4：计算实际函数值}
\[
f(x_0 + d^C) = f(-0.625, 0.781) \approx -1.940
\]

\textbf{Step 5：计算比率}
\[
\rho = \frac{f(x_0) - f(x_0 + d^C)}{q_0(0) - q_0(d^C)} = \frac{2 - (-1.940)}{3.940} = \frac{3.940}{3.940} = 1.0
\]

$\rho = 1$ 说明二次模型精确匹配实际函数（因为 $f$ 本身就是二次函数），模型预测准确，可以增大信任域半径。

\textbf{Step 6：更新}
\[
x_1 = x_0 + d^C \approx (-0.625,\ 0.781),\quad \Delta_1 > \Delta_0
\]
\end{example}

\subsection{信任域方法的Python实现}\label{sec:trust-code}

\begin{codebox}[Levenberg-Marquardt方法]
\begin{lstlisting}
import numpy as np
from scipy.optimize import least_squares

def lm_method(residuals, jac, x0, lam=0.01, nu=10, 
              eps=1e-6, max_iter=100):
    """
    Levenberg-Marquardt方法
    
    Parameters:
        residuals: 残差函数 r(x)
        jac: Jacobian函数 J(x)
        x0: 初始点
        lam: 初始正则化参数
        nu: 调节因子
        eps: 精度
        max_iter: 最大迭代次数
    
    Returns:
        x_opt: 最优解
        history: 迭代历史
    """
    x = x0.copy()
    history = [x.copy()]
    
    for k in range(max_iter):
        r = residuals(x)
        J = jac(x)
        g = J.T @ r
        
        if np.linalg.norm(g) < eps:
            break
        
        # 解 (J^T J + lam*I) d = -J^T r
        A = J.T @ J + lam * np.eye(len(x))
        d = np.linalg.solve(A, -g)
        
        # 计算比率
        rho = (np.linalg.norm(r)**2 - 
               np.linalg.norm(residuals(x + d))**2) / \
              (np.linalg.norm(r)**2 - 
               np.linalg.norm(r + J @ d)**2)
        
        if rho > 0:
            x = x + d
            history.append(x.copy())
            if rho > 0.75:
                lam = lam / nu
            elif rho < 0.25:
                lam = lam * nu
        else:
            lam = lam * nu
    
    return x, history

# 示例：非线性最小二乘拟合
# 数据
np.random.seed(42)
t = np.linspace(0, 4, 50)
y = 2.5 * np.exp(-1.3 * t) + 0.1 * np.random.randn(50)

def residuals(params):
    a, b = params
    return a * np.exp(b * t) - y

def jac(params):
    a, b = params
    J = np.zeros((len(t), 2))
    J[:, 0] = np.exp(b * t)
    J[:, 1] = a * t * np.exp(b * t)
    return J

x0 = np.array([1.0, -0.5])
x_opt, hist = lm_method(residuals, jac, x0)
print(f"拟合参数: a={x_opt[0]:.4f}, b={x_opt[1]:.4f}")
\end{lstlisting}
\end{codebox}

\begin{codebox}[信任域方法（简化版）]
\begin{lstlisting}
def trust_region(f, grad, hess, x0, Delta=1.0, eta=0.1,
                 eps=1e-6, max_iter=100):
    """
    简化信任域方法（Cauchy点策略）
    
    Parameters:
        f: 目标函数
        grad: 梯度函数
        hess: Hessian函数
        x0: 初始点
        Delta: 初始信任域半径
        eta: 接受阈值
        eps: 精度
    
    Returns:
        x_opt: 最优解
        history: 迭代历史
    """
    x = x0.copy()
    history = [x.copy()]
    Delta_max = 100.0
    
    for k in range(max_iter):
        g = grad(x)
        
        if np.linalg.norm(g) < eps:
            break
        
        H = hess(x)
        
        # 计算Cauchy点
        g_norm = np.linalg.norm(g)
        gHg = g.T @ H @ g
        
        if gHg <= 0:
            tau = 1.0
        else:
            tau = min(g_norm**3 / (Delta * gHg), 1.0)
        
        d = -tau * Delta / g_norm * g
        
        # 评估比率
        actual = f(x) - f(x + d)
        predicted = -(g.T @ d + 0.5 * d.T @ H @ d)
        
        if abs(predicted) < 1e-10:
            break
            
        rho = actual / predicted
        
        if rho > eta:
            x = x + d
            history.append(x.copy())
        
        # 更新信任域半径
        if rho < 0.25:
            Delta = 0.25 * Delta
        elif rho > 0.75 and abs(np.linalg.norm(d) - Delta) < 1e-6:
            Delta = min(2.0 * Delta, Delta_max)
    
    return x, f(x), history
\end{lstlisting}
\end{codebox}

\subsection{LM与信任域方法的比较}

\begin{table}[H]
\centering
\caption{LM与信任域方法比较}
\begin{tabular}{@{}lll@{}}
\toprule
\textbf{特性} & \textbf{LM方法} & \textbf{信任域方法} \\
\midrule
适用问题 & 最小二乘 & 一般无约束优化 \\
模型 & $J^T J + \lambda I$ & $g^T d + \frac{1}{2}d^T H d$ \\
参数 & $\lambda$（自适应） & $\Delta$（信任域半径） \\
子问题求解 & 线性方程组 & 信赖域子问题 \\
全局收敛 & 是 & 是 \\
\bottomrule
\end{tabular}
\end{table}

% 第8节：共轭方向法
\section{共轭方向法}\label{sec:conjugate}

\subsection{共轭方向的概念}\label{sec:conjugate-def}

\begin{conceptbox}[共轭方向]
对于对称正定矩阵 $A$，非零向量 $d_1, d_2, \ldots, d_k$ 称为 $A$-共轭的，如果：
\[
d_i^T A d_j = 0, \quad \forall i \neq j
\]
即它们关于 $A$ 的内积正交。
\end{conceptbox}

\begin{theorem}[共轭方向法的有限步收敛]
对于二次函数 $f(x) = \frac{1}{2}x^T A x - b^T x + c$（$A$ 正定），从任意 $x_0$ 出发，沿 $n$ 个 $A$-共轭方向精确线搜索，最多 $n$ 步收敛到最优解。
\end{theorem}

\textbf{证明}：设 $x^*$ 为最优解，$g_k = Ax_k - b$。由精确线搜索：
$g_{k+1}^T d_k = 0$

利用共轭性可证 $g_{k+1}^T d_j = 0$（$j = 0, \ldots, k$），即梯度与共轭方向正交。

$n$ 步后 $g_n$ 与 $n$ 个线性无关方向正交，故 $g_n = 0$，$x_n = x^*$。

\subsection{共轭梯度法（CG）}\label{sec:cg}

\begin{conceptbox}[共轭梯度法]
共轭梯度法的核心是通过递推自动构造\textbf{共轭方向}，无需预先给定方向组。对二次函数 $f(x) = \frac{1}{2}x^T A x - b^T x$，方向递推公式为：
\begin{align*}
d_0 &= -g_0 \\
d_{k+1} &= -g_{k+1} + \beta_k d_k
\end{align*}

\textbf{各符号含义}：
\begin{itemize}[itemsep=0.2em]
    \item $g_k = \nabla f(x_k)$：第 $k$ 步处的\textbf{梯度}（最速上升方向）
    \item $d_k$：第 $k$ 步的\textbf{搜索方向}，初始方向 $d_0 = -g_0$ 即为最速下降方向
    \item $\beta_k$：\textbf{共轭方向系数}，调节上一步方向 $d_k$ 的权重，使得新方向 $d_{k+1}$ 与 $d_k$ 关于 $A$ 共轭（即 $d_{k+1}^T A d_k = 0$）
    \item $\alpha_k$：第 $k$ 步的\textbf{最优步长}，由精确线搜索确定
\end{itemize}
\end{conceptbox}

\subsubsection{二次函数的CG算法}

\begin{algorithm}[H]
\caption{线性共轭梯度法}
\begin{algorithmic}[1]
\Require $A$（SPD），$b$，初始点 $x_0$
\State $g_0 = Ax_0 - b$，$d_0 = -g_0$
\For{$k = 0, 1, \ldots, n-1$}
    \If{$\|g_k\| = 0$} \textbf{break} \EndIf
    \State $\alpha_k = \dfrac{g_k^T g_k}{d_k^T A d_k}$
    \State $x_{k+1} = x_k + \alpha_k d_k$
    \State $g_{k+1} = g_k + \alpha_k A d_k$
    \State $\beta_k = \dfrac{g_{k+1}^T g_{k+1}}{g_k^T g_k}$
    \State $d_{k+1} = -g_{k+1} + \beta_k d_k$
\EndFor
\end{algorithmic}
\end{algorithm}

\subsubsection{CG的等价公式（$\beta_k$ 的不同形式）}

\begin{table}[H]
\centering
\caption{共轭梯度法的$\beta_k$公式}
\begin{tabular}{@{}ll@{}}
\toprule
\textbf{公式} & \textbf{表达式} \\
\midrule
Fletcher-Reeves (FR) & $\beta_k = \dfrac{g_{k+1}^T g_{k+1}}{g_k^T g_k}$ \\
Polak-Ribiere (PR) & $\beta_k = \dfrac{g_{k+1}^T(g_{k+1} - g_k)}{g_k^T g_k}$ \\
Hestenes-Stiefel (HS) & $\beta_k = \dfrac{g_{k+1}^T(g_{k+1} - g_k)}{d_k^T(g_{k+1} - g_k)}$ \\
Dai-Yuan (DY) & $\beta_k = \dfrac{g_{k+1}^T g_{k+1}}{d_k^T(g_{k+1} - g_k)}$ \\
\bottomrule
\end{tabular}
\end{table}

\begin{notebox}[PR方法]
PR方法在实际中表现最好。若 $\beta_k < 0$ 可置 $\beta_k = 0$（重启CG），保证收敛。
\end{notebox}

\subsubsection{非线性共轭梯度法}

\begin{algorithm}[H]
\caption{非线性共轭梯度法（Fletcher-Reeves）}
\begin{algorithmic}[1]
\Require 初始点 $x_0$，精度 $\varepsilon$
\State $g_0 = \nabla f(x_0)$，$d_0 = -g_0$
\For{$k = 0, 1, \ldots$}
    \If{$\|g_k\| < \varepsilon$} \textbf{break} \EndIf
    \State 线搜索求 $\alpha_k$
    \State $x_{k+1} = x_k + \alpha_k d_k$
    \State $g_{k+1} = \nabla f(x_{k+1})$
    \State $\beta_k = \dfrac{g_{k+1}^T g_{k+1}}{g_k^T g_k}$
    \State $d_{k+1} = -g_{k+1} + \beta_k d_k$
    \State 若 $k+1 \equiv 0 \pmod n$，$d_{k+1} = -g_{k+1}$（每$n$步重启）
\EndFor
\end{algorithmic}
\end{algorithm}

\begin{example}[CG法解二次函数]
$f(x) = x_1^2 + 4x_2^2$，初始点 $\begin{bmatrix} 4 \\ 1 \end{bmatrix}$。

\textbf{解：}

$A = \begin{bmatrix} 2 & 0 \\ 0 & 8 \end{bmatrix}$，$b = \begin{bmatrix} 0 \\ 0 \end{bmatrix}$

$g_0 = \begin{bmatrix} 8 \\ 8 \end{bmatrix}$，$d_0 = \begin{bmatrix} -8 \\ -8 \end{bmatrix}$

$\alpha_0 = \dfrac{g_0^T g_0}{d_0^T A d_0} = \dfrac{128}{640} = 0.2$

$x^{(1)} = \begin{bmatrix} 4 \\ 1 \end{bmatrix} + 0.2\begin{bmatrix} -8 \\ -8 \end{bmatrix} = \begin{bmatrix} 2.4 \\ -0.6 \end{bmatrix}$

$g_1 = \begin{bmatrix} 4.8 \\ -4.8 \end{bmatrix}$，$\beta_0 = \dfrac{46.08}{128} = 0.36$

$d_1 = -g_1 + \beta_0 d_0 = \begin{bmatrix} -4.8 \\ 4.8 \end{bmatrix} + 0.36\begin{bmatrix} -8 \\ -8 \end{bmatrix} = \begin{bmatrix} -7.68 \\ 1.92 \end{bmatrix}$

$\alpha_1 = \dfrac{46.08}{147.456} = 0.3125$

$x^{(2)} = \begin{bmatrix} 2.4 \\ -0.6 \end{bmatrix} + 0.3125\begin{bmatrix} -7.68 \\ 1.92 \end{bmatrix} = \begin{bmatrix} 0 \\ 0 \end{bmatrix} = x^*$！
\end{example}

\subsection{拟牛顿法}\label{sec:quasi-newton}

\begin{conceptbox}[拟牛顿法]
用矩阵 $B_k \approx H_k^{-1}$ 近似逆Hessian，满足割线条件：
\[
s_k = x_{k+1} - x_k, \quad y_k = g_{k+1} - g_k
\]
\[
B_{k+1} y_k = s_k \quad \text{（割线条件）}
\]
\end{conceptbox}

\subsubsection{DFP公式（Davidon-Fletcher-Powell）}

\begin{align}
B_{k+1}^{DFP} &= \left(I - \frac{s_k y_k^T}{y_k^T s_k}\right) B_k \left(I - \frac{y_k s_k^T}{y_k^T s_k}\right) + \frac{s_k s_k^T}{y_k^T s_k} \\
&= B_k - \frac{B_k y_k y_k^T B_k}{y_k^T B_k y_k} + \frac{s_k s_k^T}{y_k^T s_k}
\end{align}

\begin{example}[DFP方法]\label{ex:dfp}
用DFP方法求解
\[
\min\ f(x) = x_1^2 + 2x_2^2 - 4x_1 - 2x_1 x_2
\]
初始点 $x_1 = (1, 1)^T$，初始矩阵 $H_1 = I$。

\textbf{解：}

\textbf{Step 1：计算搜索方向}

梯度：$\nabla f(x) = (2x_1 - 2x_2 - 4,\ 4x_2 - 2x_1)^T$

在 $x_1 = (1,1)^T$ 处：$g_1 = (-4, 2)^T$

\[
d_1 = -H_1 g_1 = -I \begin{bmatrix} -4 \\ 2 \end{bmatrix} = \begin{bmatrix} 4 \\ -2 \end{bmatrix}
\]

\textbf{Step 2：精确线搜索}

令 $x_1 + \alpha_1 d_1 = (1+4\alpha_1,\ 1-2\alpha_1)^T$，代入 $f$ 求 $\alpha_1$：
\[
\phi(\alpha_1) = f(1+4\alpha_1,\ 1-2\alpha_1) = 40\alpha_1^2 - 24\alpha_1 - 1
\]
$\phi'(\alpha_1) = 80\alpha_1 - 24 = 0$，得 $\alpha_1 = \dfrac{3}{10}$。

\[
x_2 = \begin{bmatrix} 1 \\ 1 \end{bmatrix} + \frac{3}{10}\begin{bmatrix} 4 \\ -2 \end{bmatrix} = \begin{bmatrix} 11/5 \\ 4/5 \end{bmatrix}
\]

\textbf{Step 3：计算 $s_1, y_1$}

\[
g_2 = \nabla f\!\left(\frac{11}{5}, \frac{4}{5}\right) = \left(\frac{6}{5},\, -\frac{6}{5}\right)^T
\]
\[
s_1 = x_2 - x_1 = \left(\frac{6}{5},\, -\frac{6}{5}\right)^T, \quad
y_1 = g_2 - g_1 = \left(\frac{26}{5},\, -\frac{16}{5}\right)^T
\]

\textbf{Step 4：DFP更新}

由 $H_1 = I$，按DFP公式：
\[
H_2 = H_1 - \frac{H_1 y_1 y_1^T H_1}{y_1^T H_1 y_1} + \frac{s_1 s_1^T}{y_1^T s_1}
\]

逐项计算：
\[
y_1^T H_1 y_1 = y_1^T y_1 = \frac{676 + 256}{25} = \frac{932}{25}
\]
\[
\frac{H_1 y_1 y_1^T H_1}{y_1^T H_1 y_1} = \frac{1}{932}\begin{bmatrix} 676 & -416 \\ -416 & 256 \end{bmatrix}
= \begin{bmatrix} 169/233 & -104/233 \\ -104/233 & 64/233 \end{bmatrix}
\]
\[
y_1^T s_1 = \frac{156 + 96}{25} = \frac{252}{25}, \quad
\frac{s_1 s_1^T}{y_1^T s_1} = \frac{1}{252}\begin{bmatrix} 36 & -36 \\ -36 & 36 \end{bmatrix}
= \begin{bmatrix} 1/7 & -1/7 \\ -1/7 & 1/7 \end{bmatrix}
\]

合并得：
\[
H_2 = \begin{bmatrix} 1 - \dfrac{169}{233} + \dfrac{1}{7} & \dfrac{104}{233} - \dfrac{1}{7} \\[8pt]
\dfrac{104}{233} - \dfrac{1}{7} & 1 - \dfrac{64}{233} + \dfrac{1}{7} \end{bmatrix}
= \begin{bmatrix} \dfrac{681}{1631} & \dfrac{495}{1631} \\[6pt] \dfrac{495}{1631} & \dfrac{1416}{1631} \end{bmatrix}
\approx \begin{bmatrix} 0.418 & 0.304 \\ 0.304 & 0.868 \end{bmatrix}
\]

\textbf{Step 5：新的搜索方向}
\[
d_2 = -H_2 g_2 = -\begin{bmatrix} 681/1631 & 495/1631 \\ 495/1631 & 1416/1631 \end{bmatrix}
\begin{bmatrix} 6/5 \\ -6/5 \end{bmatrix}
= \begin{bmatrix} -\dfrac{1116}{8155} \\[6pt] \dfrac{5526}{8155} \end{bmatrix}
\approx \begin{bmatrix} -0.137 \\ 0.678 \end{bmatrix}
\]
\end{example}

\begin{algorithm}[H]
\caption{DFP方法}
\begin{algorithmic}[1]
\Require 初始点 $x_0$，初始矩阵 $H_0$（通常 $H_0 = I$）
\State $k \gets 0$
\While{$\|\nabla f(x_k)\| \geq \varepsilon$}
    \State $d_k = -H_k \nabla f(x_k)$
    \State 线搜索求 $\alpha_k$（满足Wolfe条件）
    \State $s_k = \alpha_k d_k$，$x_{k+1} = x_k + s_k$
    \State $y_k = \nabla f(x_{k+1}) - \nabla f(x_k)$
    \If{$y_k^T s_k > 0$}
        \State $H_{k+1} = H_k - \dfrac{H_k y_k y_k^T H_k}{y_k^T H_k y_k} + \dfrac{s_k s_k^T}{y_k^T s_k}$
    \Else
        \State $H_{k+1} = H_k$（或重启）
    \EndIf
    \State $k \gets k + 1$
\EndWhile
\end{algorithmic}
\end{algorithm}

\begin{remark}
DFP 方法维护的是逆 Hessian 近似 $H_k \approx (\nabla^2 f(x_k))^{-1}$，因此搜索方向直接由 $d_k = -H_k g_k$ 给出，无需像 BFGS 那样求解线性方程组。但 DFP 对线搜索精度要求更高，数值稳定性不如 BFGS，实践中 BFGS 更为常用。
\end{remark}

\subsubsection{BFGS公式（Broyden-Fletcher-Goldfarb-Shanno）}

\begin{align}
B_{k+1}^{BFGS} &= B_k + \frac{(s_k - B_k y_k)(s_k - B_k y_k)^T}{y_k^T(s_k - B_k y_k)} \quad \text{(秩一修正)} \\
&= \left(I - \frac{s_k y_k^T}{y_k^T s_k}\right) B_k \left(I - \frac{y_k s_k^T}{y_k^T s_k}\right) + \frac{s_k s_k^T}{y_k^T s_k}
\end{align}

或等价地对 $H$ 的更新：
\[
H_{k+1}^{BFGS} = H_k - \frac{H_k s_k s_k^T H_k}{s_k^T H_k s_k} + \frac{y_k y_k^T}{y_k^T s_k}
\]

\begin{algorithm}[H]
\caption{BFGS方法}
\begin{algorithmic}[1]
\Require 初始点 $x_0$，初始矩阵 $B_0$（通常 $B_0 = I$）
\State $k \gets 0$
\While{$\|\nabla f(x_k)\| \geq \varepsilon$}
    \State $d_k = -B_k \nabla f(x_k)$
    \State 线搜索求 $\alpha_k$（满足Wolfe条件）
    \State $s_k = \alpha_k d_k$，$x_{k+1} = x_k + s_k$
    \State $y_k = \nabla f(x_{k+1}) - \nabla f(x_k)$
    \If{$y_k^T s_k > 0$}
        \State 用BFGS公式更新 $B_{k+1}$
    \Else
        \State $B_{k+1} = B_k$（或重启）
    \EndIf
    \State $k \gets k + 1$
\EndWhile
\end{algorithmic}
\end{algorithm}

\begin{theorem}[BFGS的收敛性]
若 $f$ 为凸函数且 $L$-光滑，采用Wolfe条件线搜索，则BFGS全局收敛到最优解。若 $f$ 还强凸，则超线性收敛。
\end{theorem}

\begin{example}[BFGS求解]
用BFGS方法求解 $\min\, f(x) = x_1^2 + 4x_2^2$，初始点 $x_0 = (4, 1)^T$。

\textbf{解：}

取 $B_0 = I$。计算 $g_0 = \nabla f(x_0) = (8, 8)^T$。

搜索方向 $d_0 = -B_0 g_0 = (-8, -8)^T$，精确线搜索得 $\alpha_0 = 0.2$。

$x_1 = (2.4, -0.6)$，$g_1 = (4.8, -4.8)$

$s_0 = (-1.6, -1.6)$，$y_0 = g_1 - g_0 = (-3.2, -12.8)$

计算BFGS更新得 $B_1 = \begin{bmatrix} 0.5 & 0 \\ 0 & 0.125 \end{bmatrix} = A^{-1}$

$d_1 = -B_1 g_1 = (-2.4, 0.6)$，$x_2 = x_1 + d_1 = (0, 0) = x^*$！
\end{example}

\subsection{共轭梯度法与拟牛顿法的比较}

\begin{table}[H]
\centering
\caption{共轭梯度法与拟牛顿法比较}
\begin{tabular}{@{}lll@{}}
\toprule
\textbf{特性} & \textbf{CG} & \textbf{BFGS} \\
\midrule
内存需求 & $O(n)$ & $O(n^2)$ \\
每次迭代计算量 & 低 & 中等 \\
收敛速度 & 线性/超线性 & 超线性 \\
适合问题规模 & 大规模 & 中大规模 \\
重启策略 & 需要 & 不需要 \\
全局收敛 & 需配合线搜索 & 配合Wolfe条件 \\
\bottomrule
\end{tabular}
\end{table}

\subsection{Python代码实现}\label{sec:conjugate-code}

\begin{codebox}[共轭梯度法]
\begin{lstlisting}
import numpy as np
from scipy.optimize import minimize_scalar

def conjugate_gradient(f, grad, x0, method='FR', eps=1e-6,
                       max_iter=1000):
    """
    非线性共轭梯度法
    
    Parameters:
        f: 目标函数
        grad: 梯度函数
        x0: 初始点
        method: 'FR', 'PR', 'HS'
        eps: 精度
        max_iter: 最大迭代次数
    
    Returns:
        x_opt: 最优解
        history: 迭代历史
    """
    x = x0.copy()
    n = len(x)
    g = grad(x)
    d = -g.copy()
    history = [x.copy()]
    
    k = 0
    for _ in range(max_iter):
        if np.linalg.norm(g) < eps:
            break
        
        # 线搜索
        def phi(alpha):
            return f(x + alpha * d)
        res = minimize_scalar(phi)
        alpha = res.x
        
        x = x + alpha * d
        history.append(x.copy())
        
        g_new = grad(x)
        
        # 计算 beta
        if method == 'FR':
            beta = np.dot(g_new, g_new) / np.dot(g, g)
        elif method == 'PR':
            beta = np.dot(g_new, g_new - g) / np.dot(g, g)
            beta = max(0, beta)  # PR+ 保证收敛
        elif method == 'HS':
            dy = g_new - g
            beta = np.dot(g_new, dy) / np.dot(d, dy)
        
        # 重启
        k += 1
        if k % n == 0:
            d = -g_new
        else:
            d = -g_new + beta * d
        
        g = g_new
    
    return x, f(x), history

# 示例
x0 = np.array([4.0, 1.0])
x_opt, f_opt, hist = conjugate_gradient(
    lambda x: x[0]**2 + 4*x[1]**2,
    lambda x: np.array([2*x[0], 8*x[1]]),
    x0
)
print(f"最优解: {x_opt}")
print(f"最优值: {f_opt:.8f}")
\end{lstlisting}
\end{codebox}

\begin{codebox}[BFGS方法]
\begin{lstlisting}
def bfgs(f, grad, x0, eps=1e-6, max_iter=1000):
    """
    BFGS拟牛顿法
    
    Parameters:
        f: 目标函数
        grad: 梯度函数
        x0: 初始点
        eps: 精度
        max_iter: 最大迭代次数
    
    Returns:
        x_opt: 最优解
        history: 迭代历史
    """
    n = len(x0)
    x = x0.copy()
    B = np.eye(n)  # 初始Hessian逆近似
    g = grad(x)
    history = [x.copy()]
    
    for _ in range(max_iter):
        if np.linalg.norm(g) < eps:
            break
        
        # 搜索方向
        d = -B @ g
        
        # Wolfe条件线搜索
        alpha = 1.0
        c1, c2 = 1e-4, 0.9
        fx = f(x)
        
        for _ in range(20):
            x_new = x + alpha * d
            g_new = grad(x_new)
            
            # 充分下降
            if f(x_new) <= fx + c1 * alpha * np.dot(g, d):
                # 曲率条件
                if np.dot(g_new, d) >= c2 * np.dot(g, d):
                    break
            alpha *= 0.5
        
        s = alpha * d
        x = x + s
        g_new = grad(x)
        y = g_new - g
        
        # BFGS更新
        if np.dot(y, s) > 1e-10:  # 保证正定性
            rho = 1.0 / np.dot(y, s)
            I = np.eye(n)
            B = (I - rho * np.outer(s, y)) @ B @ \
                (I - rho * np.outer(y, s)) + \
                rho * np.outer(s, s)
        
        g = g_new
        history.append(x.copy())
    
    return x, f(x), history

# 示例：Rosenbrock函数
x0 = np.array([-1.0, 2.0])
x_opt, f_opt, hist = bfgs(rosenbrock, rosen_grad, x0)
print(f"最优解: {x_opt}")
print(f"最优值: {f_opt:.10f}")
print(f"迭代次数: {len(hist)}")
\end{lstlisting}
\end{codebox}

% 第9节：次梯度优化方法
\section{次梯度优化方法}\label{sec:subgradient}

\subsection{背景与动机}

许多优化问题涉及非光滑目标函数（如L1正则化、hinge loss等），梯度不存在，需要推广导数的概念。

\subsection{次梯度与次微分}\label{sec:subgradient-def}

\begin{definition}[次梯度]
设 $f: \mathbb{R}^n \to \mathbb{R}$ 为凸函数。向量 $g \in \mathbb{R}^n$ 称为 $f$ 在 $x$ 处的次梯度，如果：
\[
f(y) \geq f(x) + g^T(y - x), \quad \forall y \in \mathbb{R}^n
\]
几何意义：$g$ 定义了 $f$ 在 $x$ 处的一个支撑超平面。
\end{definition}

\begin{definition}[次微分]
$f$ 在 $x$ 处的所有次梯度的集合称为次微分，记为 $\partial f(x)$：
\[
\partial f(x) = \{g : f(y) \geq f(x) + g^T(y - x), \forall y\}
\]
\end{definition}

\begin{theorem}[次微分的性质]
\begin{enumerate}
    \item $\partial f(x)$ 是非空紧凸集
    \item 若 $f$ 在 $x$ 处可微，则 $\partial f(x) = \{\nabla f(x)\}$
    \item $x^*$ 是全局最优 $\Leftrightarrow 0 \in \partial f(x^*)$
    \item（正齐次性）$\partial(\alpha f)(x) = \alpha \partial f(x)$（$\alpha > 0$）
    \item（加法规则）$\partial(f_1 + f_2)(x) \supseteq \partial f_1(x) + \partial f_2(x)$
\end{enumerate}
\end{theorem}

\begin{figure}[H]
\centering
\begin{tikzpicture}[scale=0.9]
    % 坐标轴
    \draw[->,thick] (-3.2,0) -- (3.2,0) node[right] {$x$};
    \draw[->,thick] (0,-0.3) -- (0,3.5) node[above] {$f(x)$};
    \node[below left] at (0,0) {$O$};
    % 绝对值函数图像 f(x) = |x|
    \draw[very thick,blue] (-3,3) -- (0,0) -- (3,3);
    \node[blue,right] at (3.1,2.9) {$f(x) = |x|$};
    % 在 x=1.5 处的次梯度（唯一，g = 1）
    \fill[red] (1.5,1.5) circle (2.5pt) node[above right,scale=0.8] {$(1.5, 1.5)$};
    \draw[thick,red,dashed] (-0.5,-0.5) -- (3,3);
    \node[red,scale=0.75] at (3.3,2.7) {\small $g=1$};
    % 在 x=-1 处的次梯度（唯一，g = -1）
    \fill[red] (-1,1) circle (2.5pt) node[above left,scale=0.8] {$(-1, 1)$};
    \draw[thick,green!50!black,dashed] (-3,3) -- (1,-1);
    \node[green!50!black,scale=0.75] at (1.3,-0.7) {\small $g=-1$};
    % 在 x=0 处的次梯度锥（整个区间 [-1, 1]）
    \fill[orange] (0,0) circle (3pt);
    \draw[thick,purple,dashed] (-2.5,2.5) -- (2.5,-2.5);
    \node[purple,scale=0.75] at (2.8,-1.5) {\small $g=-1$};
    \draw[thick,purple,dashed] (-2.5,-2.5) -- (2.5,2.5);
    \node[purple,scale=0.75] at (2.8,1.5) {\small $g=1$};
    \node[purple,scale=0.75] at (1.8,0.3) {\small 次微分 $[-1,1]$};
\end{tikzpicture}
\caption{次梯度与次微分的几何解释：$f(x)=|x|$ 在 $x \neq 0$ 处有唯一次梯度 $g = \operatorname{sign}(x)$，在 $x=0$ 处次微分为区间 $[-1, 1]$}
\label{fig:subgradient-geom}
\end{figure}

\subsection{常见函数的次微分}\label{sec:common-subgradients}

\begin{example}[绝对值函数的次微分]
$f(x) = |x|$：

\textbf{解：}
\[
\partial f(x) = \begin{cases} \{-1\} & x < 0 \\ \{-1, 1\} & x = 0 \\ \{+1\} & x > 0 \end{cases}
\]
\end{example}

\begin{example}[欧几里得范数的次微分]
$f(x) = \|x\|$：

\textbf{解：}
\[
\partial f(x) = \begin{cases} \left\{\frac{x}{\|x\|}\right\} & x \neq 0 \\ \{g : \|g\| \leq 1\} & x = 0 \end{cases}
\]
\end{example}

\begin{example}[最大值函数的次微分]
$f(x) = \max_{i=1,\ldots,m} f_i(x)$：

\textbf{解：}
\[
\partial f(x) = \text{conv}\left(\bigcup_{i \in I(x)} \partial f_i(x)\right)
\]
其中 $I(x) = \{i : f_i(x) = f(x)\}$ 为活跃指标集。
\end{example}

\subsection{次梯度方法}\label{sec:subgradient-method}

\begin{conceptbox}[次梯度方法]
直接推广梯度下降：
\[
x_{k+1} = x_k - \alpha_k g_k, \quad g_k \in \partial f(x_k)
\]
与梯度下降不同，$-\!g_k$ 不一定是下降方向。
\end{conceptbox}

\begin{theorem}[次梯度方法的收敛性]
设 $f$ 为凸函数，$\|g_k\| \leq G$，$\|x_0 - x^*\| \leq R$。

\textbf{固定步长} $\alpha_k = \alpha$：
\[
\min_{0 \leq i < k} f(x_i) - f^* \leq \frac{R^2 + G^2 \alpha^2 k}{2\alpha k} \to \frac{G^2 \alpha}{2}
\]
收敛到 $O(\alpha)$ 邻域内。

\textbf{递减步长} $\alpha_k = \frac{R}{G\sqrt{k}}$：
\[
\min_{0 \leq i < k} f(x_i) - f^* \leq \frac{RG}{\sqrt{k}} = O\left(\frac{1}{\sqrt{k}}\right)
\]
\end{theorem}

\begin{notebox}[次梯度方法的特点]
\begin{itemize}
    \item 收敛速率 $O(1/\sqrt{k})$，比梯度下降慢
    \item 步长选择至关重要（递减步长）
    \item 函数值不单调下降
    \item 适合大规模问题（每次迭代简单）
\end{itemize}
\end{notebox}

\subsubsection{投影次梯度方法}

对于约束问题 $\min_{x \in C} f(x)$：
\[
x_{k+1} = P_C(x_k - \alpha_k g_k)
\]
其中 $P_C(x) = \arg\min_{y \in C} \|y - x\|$ 为到集合 $C$ 的投影。

\subsection{近端梯度方法（Proximal Gradient）}\label{sec:proximal}

\begin{conceptbox}[近端算子]
对于函数 $h$ 和参数 $\lambda > 0$，近端算子定义为：
\[
\text{prox}_{\lambda h}(x) = \arg\min_u \left\{h(u) + \frac{1}{2\lambda}\|u - x\|^2\right\}
\]
\end{conceptbox}

对于复合优化问题：$\min f(x) = g(x) + h(x)$

其中 $g$ 光滑，$h$ 非光滑但简单（近端可计算）。

\textbf{近端梯度法（ISTA）}：
\[
x_{k+1} = \text{prox}_{\alpha_k h}(x_k - \alpha_k \nabla g(x_k))
\]

\begin{example}[LASSO问题]
$\min \frac{1}{2}\|Ax - b\|^2 + \lambda\|x\|_1$

\textbf{解：}

$g(x) = \frac{1}{2}\|Ax - b\|^2$，$h(x) = \lambda\|x\|_1$

近端算子（软阈值）：
\[
[\text{prox}_{\lambda h}(x)]_i = \text{sign}(x_i) \max(|x_i| - \lambda, 0)
\]

迭代公式（ISTA）：
\[
x_{k+1} = S_{\alpha \lambda}(x_k - \alpha A^T(Ax_k - b))
\]
其中 $S$ 为软阈值算子。
\end{example}

\subsection{加速近端梯度法（FISTA）}\label{sec:fista}

\begin{algorithm}[H]
\caption{FISTA}
\begin{algorithmic}[1]
\Require 初始点 $x_0 = y_1$，步长 $\alpha \leq 1/L$
\State $t_1 = 1$
\For{$k = 1, 2, \ldots$}
    \State $x_k = \text{prox}_{\alpha h}(y_k - \alpha \nabla g(y_k))$
    \State $t_{k+1} = \frac{1 + \sqrt{1 + 4t_k^2}}{2}$
    \State $y_{k+1} = x_k + \frac{t_k - 1}{t_{k+1}}(x_k - x_{k-1})$
\EndFor
\end{algorithmic}
\end{algorithm}

\begin{theorem}[FISTA收敛率]
FISTA的收敛率为 $O(1/k^2)$：
\[
f(x_k) - f^* \leq \frac{2L\|x_0 - x^*\|^2}{(k+1)^2}
\]
\end{theorem}

\subsection{次梯度方法的Python实现}\label{sec:subgradient-code}

\begin{codebox}[次梯度方法]
\begin{lstlisting}
import numpy as np

def subgradient_method(f, subgrad, x0, step_rule='diminishing',
                       alpha=0.01, eps=1e-6, max_iter=1000):
    """
    次梯度方法
    
    Parameters:
        f: 目标函数
        subgrad: 次梯度函数
        x0: 初始点
        step_rule: 'fixed', 'diminishing', 'polyak'
        alpha: 步长参数
        eps: 精度
        max_iter: 最大迭代次数
    
    Returns:
        x_best: 最优解
        history: 函数值历史
    """
    x = x0.copy()
    f_best = f(x)
    x_best = x.copy()
    history = [f_best]
    
    for k in range(1, max_iter + 1):
        g = subgrad(x)
        
        if np.linalg.norm(g) < eps:
            break
        
        # 步长选择
        if step_rule == 'fixed':
            step = alpha
        elif step_rule == 'diminishing':
            step = alpha / np.sqrt(k)
        elif step_rule == 'polyak':
            # 需要知道最优值 f*
            step = (f(x) - 0) / (np.linalg.norm(g)**2 + 1e-10)
        
        x = x - step * g
        f_val = f(x)
        history.append(f_val)
        
        # 记录最优
        if f_val < f_best:
            f_best = f_val
            x_best = x.copy()
    
    return x_best, history

# 示例：绝对值函数求和
f = lambda x: np.sum(np.abs(x))
subgrad = lambda x: np.sign(x)  # 在零点取0

x0 = np.array([5.0, -3.0, 2.0])
x_best, hist = subgradient_method(f, subgrad, x0, 
                                  step_rule='diminishing',
                                  alpha=1.0)
print(f"最优解: {x_best}")
print(f"最优值: {f(x_best):.6f}")
\end{lstlisting}
\end{codebox}

\begin{codebox}[近端梯度法（ISTA/FISTA）]
\begin{lstlisting}
def ista(g, grad_g, prox_h, x0, alpha=0.01, eps=1e-6,
         max_iter=1000):
    """
    近端梯度法（ISTA）
    
    Parameters:
        g: 光滑部分函数
        grad_g: 光滑部分梯度
        prox_h: 近端算子
        x0: 初始点
        alpha: 步长
        eps: 精度
    
    Returns:
        x_opt: 最优解
        history: 迭代历史
    """
    x = x0.copy()
    history = [g(x) + 0]  # 需要外部提供完整f值
    
    for _ in range(max_iter):
        x_new = prox_h(x - alpha * grad_g(x), alpha)
        
        if np.linalg.norm(x_new - x) < eps:
            break
        
        x = x_new
        history.append(x.copy())
    
    return x, history

def fista(g, grad_g, prox_h, x0, alpha=0.01, max_iter=1000):
    """
    加速近端梯度法（FISTA）
    """
    x = x0.copy()
    y = x0.copy()
    t = 1.0
    history = [x.copy()]
    
    for _ in range(max_iter):
        x_prev = x.copy()
        
        x = prox_h(y - alpha * grad_g(y), alpha)
        
        t_new = (1 + np.sqrt(1 + 4 * t**2)) / 2
        y = x + ((t - 1) / t_new) * (x - x_prev)
        
        t = t_new
        history.append(x.copy())
    
    return x, history

# 软阈值算子（L1近端算子）
def soft_threshold(x, lam):
    return np.sign(x) * np.maximum(np.abs(x) - lam, 0)

# LASSO: min 0.5*||Ax-b||^2 + lambda*||x||_1
def solve_lasso(A, b, lam, alpha=None):
    n = A.shape[1]
    if alpha is None:
        alpha = 1.0 / np.linalg.norm(A.T @ A, 2)
    
    grad_g = lambda x: A.T @ (A @ x - b)
    prox_h = lambda x, alpha: soft_threshold(x, lam * alpha)
    
    x0 = np.zeros(n)
    x_opt, hist = fista(None, grad_g, prox_h, x0, alpha,
                        max_iter=1000)
    return x_opt, hist
\end{lstlisting}
\end{codebox}

\subsection{复杂度下界与加速方法}

\begin{theorem}[一阶方法的复杂度下界]
对于 $L$-光滑的凸函数，任何一阶方法的收敛速率下界为 $O(1/k^2)$。

对于 $\mu$-强凸且 $L$-光滑的函数，下界为：
\[
f(x_k) - f^* \geq \left(\frac{\sqrt{\kappa} - 1}{\sqrt{\kappa} + 1}\right)^{2k} (f(x_0) - f^*)
\]
其中 $\kappa = L/\mu$ 为条件数。
\end{theorem}

\begin{notebox}[Nesterov加速梯度法]
对于梯度下降 $O(1/k)$ 的收敛率，Nesterov加速可达 $O(1/k^2)$：

\textbf{加速梯度法}：
\begin{align*}
y_k &= x_k + \beta_k(x_k - x_{k-1}) \\
x_{k+1} &= y_k - \alpha_k \nabla f(y_k)
\end{align*}
其中 $\beta_k = \frac{k-1}{k+2}$（凸情形）或 $\beta_k = \frac{\sqrt{\kappa} - 1}{\sqrt{\kappa} + 1}$（强凸情形）。
\end{notebox}

% 第10节：总结
\section{本章小结}\label{sec:ch5-summary}

\subsection{核心思想}

\begin{conceptbox}[多维无约束优化的两大支柱]
\begin{enumerate}[label=\arabic*., leftmargin=2em]
    \item \textbf{下降迭代框架}：$x_{k+1} = x_k + \alpha_k d_k$，其中 $d_k$ 为下降方向，$\alpha_k$ 为步长（由第~\ref{chap:linesearch}~章的线搜索方法确定）。
    \item \textbf{以信息代价换取收敛速度}：仅函数值（零阶）$\rightarrow$ 梯度（一阶）$\rightarrow$ Hessian（二阶），信息越多，收敛越快。
\end{enumerate}
\end{conceptbox}

\subsection{方法体系总览}

\begin{table}[H]
\centering
\small
\caption{无约束最优化方法总览（步长搜索统一使用第~\ref{chap:linesearch}~章的线搜索方法）}
\setlength{\tabcolsep}{4pt}
\begin{tabular}{@{}c>{\raggedright\arraybackslash}p{3cm}>{\raggedright\arraybackslash}p{2.2cm}>{\raggedright\arraybackslash}p{4cm}@{}}
\toprule
\textbf{信息类型} & \textbf{代表方法} & \textbf{收敛速度} & \textbf{适用场景} \\
\midrule
\multirow{3}{*}{\textbf{零阶}} & 坐标轮换法 & 线性 & 小规模 \\
 & Hooke-Jeeves & 线性 & 模式搜索 \\
 & Rosenbrock & 超线性 & 峡谷函数 \\
\midrule
\multirow{3}{*}{\textbf{一阶}} & 最速下降法 & $O(1/k)$ & 大规模 \\
 & CG & 超线性 & 大规模稀疏 \\
 & SGD & $O(1/\sqrt{k})$ & 机器学习 \\
\midrule
\multirow{3}{*}{\textbf{二阶}} & 牛顿法 & 局部二次 & 中小规模 \\
 & Levenberg-Marquardt & 超线性 & 最小二乘 \\
 & 信任域方法 & 超线性 & 非凸问题 \\
\midrule
\textbf{拟牛顿} & DFP、BFGS & 超线性 & 中大规模 \\
\midrule
\multirow{3}{*}{\textbf{非光滑}} & 次梯度方法 & $O(1/\sqrt{k})$ & 凸非光滑 \\
 & 近端梯度法 & $O(1/k)$ & 复合优化 \\
 & FISTA & $O(1/k^2)$ & 复合优化 \\
\bottomrule
\end{tabular}
\end{table}

\subsection{方法选择决策树}

\begin{conceptbox}[如何选择多维无约束优化方法]
\begin{enumerate}[label=\arabic*., leftmargin=2em]
    \item \textbf{函数是否光滑？}
    \begin{itemize}
        \item 非光滑凸 $\rightarrow$ 近端梯度法 / FISTA
        \item 非光滑且导数不可得 $\rightarrow$ 坐标轮换法 / Hooke-Jeeves
        \item 光滑 $\rightarrow$ 继续判断
    \end{itemize}
    \item \textbf{能否计算Hessian？}
    \begin{itemize}
        \item 能且 $n < 1000$ $\rightarrow$ 牛顿法（二次收敛）或信任域方法（全局收敛）
        \item 不能 $\rightarrow$ 继续判断
    \end{itemize}
    \item \textbf{问题规模？}
    \begin{itemize}
        \item $n > 10^5$（大规模）$\rightarrow$ 共轭梯度法 / SGD
        \item $1000 < n < 10^5$（中规模）$\rightarrow$ BFGS（实用首选）
        \item $n < 1000$（小规模）$\rightarrow$ BFGS 或 牛顿法
    \end{itemize}
    \item \textbf{特殊结构？}
    \begin{itemize}
        \item 最小二乘问题 $\rightarrow$ LM方法
        \item 复合优化（$f = g + h$，$g$光滑$h$非光滑）$\rightarrow$ 近端梯度 / FISTA
        \item 大数据/随机采样 $\rightarrow$ SGD及其变体
    \end{itemize}
\end{enumerate}
\end{conceptbox}

\subsection{关键概念回顾}

\begin{enumerate}[label=\arabic*., leftmargin=2em]
    \item \textbf{下降方向}：$d$ 满足 $\nabla f(x)^T d < 0$；最速下降方向 $d = -\nabla f(x)$
    
    \item \textbf{共轭性}：$d_i^T A d_j = 0$（$i \neq j$），保证二次函数 $n$ 步内收敛；CG法通过递推构造共轭方向
    
    \item \textbf{割线条件}：$B_{k+1} s_k = y_k$（其中 $s_k = x_{k+1}-x_k$，$y_k = \nabla f_{k+1}-\nabla f_k$），拟Newton法通过此条件更新Hessian近似
    
    \item \textbf{信任域 vs 线搜索}：信任域先限制步长范围再优化方向，全局收敛性更好；线搜索先确定方向再找步长
    
    \item \textbf{次梯度}：$g \in \partial f(x)$，满足 $f(y) \geq f(x) + g^T(y-x)$，推广梯度到非光滑凸函数
    
    \item \textbf{近端算子}：$\text{prox}_{\lambda h}(x) = \arg\min_y \{h(y) + \frac{1}{2\lambda}\|y-x\|^2\}$，将非光滑项的极小化显式求解
    
    \item \textbf{加速技术}：FISTA通过Nesterov动量将近端梯度法从 $O(1/k)$ 加速到 $O(1/k^2)$
\end{enumerate}

\subsection{收敛速度总结}

\begin{table}[H]
\centering
\caption{各方法收敛速度对照}
\begin{tabular}{@{}lccc@{}}
\toprule
\textbf{方法} & \textbf{一般凸} & \textbf{强凸} & \textbf{局部（光滑）} \\
\midrule
坐标轮换法 & $-$ & $-$ & 线性 \\
最速下降法 & $O(1/k)$ & $O((1-\mu/L)^k)$ & 线性 \\
共轭梯度法 & $-$ & $-$ & 超线性 \\
牛顿法 & $-$ & $-$ & 二次 \\
BFGS & $-$ & 超线性 & 超线性 \\
SGD & $O(1/\sqrt{k})$ & $O(1/k)$ & $-$ \\
次梯度方法 & $O(1/\sqrt{k})$ & $-$ & $-$ \\
近端梯度法 & $O(1/k)$ & 线性 & $-$ \\
FISTA & $O(1/k^2)$ & $-$ & $-$ \\
\bottomrule
\end{tabular}
\end{table}

\subsection{关键公式汇总}

\begin{table}[H]
\centering
\caption{核心公式速查}
\small
\begin{tabular}{@{}lp{9.5cm}@{}}
\toprule
\textbf{概念/方法} & \textbf{公式} \\
\midrule
最速下降方向 & $d_k = -\nabla f(x_k)$ \\
Newton方向 & $d_k = -[\nabla^2 f(x_k)]^{-1} \nabla f(x_k)$ \\
CG（Fletcher-Reeves） & $d_k = -g_k + \beta_k d_{k-1}$，$\beta_k = \|g_k\|^2/\|g_{k-1}\|^2$ \\
CG（Polak-Ribi\`ere） & $\beta_k = g_k^T(g_k-g_{k-1})/\|g_{k-1}\|^2$ \\
DFP更新 & $B_{k+1} = B_k - \frac{B_k y_k y_k^T B_k}{y_k^T B_k y_k} + \frac{s_k s_k^T}{y_k^T s_k}$ \\
BFGS更新 & $H_{k+1} = V_k^T H_k V_k + \rho_k s_k s_k^T$，$\rho_k = 1/(y_k^T s_k)$ \\
LM方法 & $(J^T J + \lambda I) d = -J^T r$ \\
信任域子问题 & $\min_{\|d\|\leq\Delta} m_k(d) = f_k + g_k^T d + \frac{1}{2}d^T B_k d$ \\
次梯度迭代 & $x_{k+1} = x_k - \alpha_k g_k$，$g_k \in \partial f(x_k)$ \\
近端梯度法 & $x_{k+1} = \text{prox}_{\alpha_k h}(x_k - \alpha_k \nabla g(x_k))$ \\
FISTA & $x_k = \text{prox}_{\alpha h}(y_k - \alpha \nabla g(y_k))$ \\
 & $y_{k+1} = x_k + \frac{k-1}{k+2}(x_k - x_{k-1})$ \\
\bottomrule
\end{tabular}
\end{table}

\subsection{学习建议}

\begin{notebox}[学习路径]
\begin{enumerate}[label=\arabic*., leftmargin=2em]
    \item 掌握最速下降法的收敛性分析和步长选择（理解线搜索在多维优化中的作用）
    \item 理解牛顿法的二次收敛性和局限性（Hessian计算代价、不正定问题）
    \item 掌握BFGS和CG作为实用的大规模方法（拟Newton近似、共轭方向构造）
    \item 了解信任域方法如何保证全局收敛（与线搜索策略的对比）
    \item 了解次梯度方法处理非光滑问题（凸分析与近端算子）
\end{enumerate}
\end{notebox}

\subsection{关键术语中英对照}

\begin{table}[H]
\centering
\begin{tabular}{@{}ll@{}}
\toprule
\textbf{中文} & \textbf{英文} \\
\midrule
最速下降法 & Steepest Descent / Gradient Descent \\
牛顿法 & Newton's Method \\
共轭梯度法 & Conjugate Gradient (CG) \\
拟牛顿法 & Quasi-Newton Method \\
DFP公式 & Davidon-Fletcher-Powell Formula \\
BFGS公式 & Broyden-Fletcher-Goldfarb-Shanno Formula \\
信任域方法 & Trust Region Method \\
Levenberg-Marquardt & Levenberg-Marquardt Algorithm \\
次梯度 & Subgradient \\
次微分 & Subdifferential \\
近端算子 & Proximal Operator \\
近端梯度法 & Proximal Gradient Method \\
FISTA & Fast Iterative Shrinkage-Thresholding Algorithm \\
坐标轮换法 & Cyclic Coordinate Method \\
Hooke-Jeeves方法 & Hooke-Jeeves Pattern Search \\
Rosenbrock方法 & Rosenbrock's Method \\
下降方向 & Descent Direction \\
共轭方向 & Conjugate Direction \\
割线条件 & Secant Condition \\
\bottomrule
\end{tabular}
\end{table}

\newpage

% ==================== 第六章 约束最优化理论 ====================
